\chapter{Marco de Referencia}

\section{Marco Teórico}
Para comprender el desarrollo de un módulo de ERP destinado a la automatización de la asignación de tareas en un taller industrial, es esencial establecer una base teórica sólida sobre los conceptos clave involucrados. Este proyecto se enfoca en la aplicación de algoritmos genéticos como una técnica de inteligencia artificial que permite la automatización de procesos de toma de decisiones en entornos industriales. Los algoritmos genéticos permiten analizar múltiples combinaciones posibles para la asignación de recursos y actividades, facilitando la gestión eficiente de órdenes de trabajo en un ERP.

Dentro del campo de la ingeniería industrial, se han desarrollado diversas metodologías para mejorar la productividad y la gestión de recursos, muchas de ellas basadas en principios de automatización y sistemas inteligentes. En este proyecto, se desarrollará un módulo de ERP que empleará algoritmos genéticos para la automatización de la asignación de tareas, permitiendo reducir la intervención manual y mejorar la eficiencia operativa en talleres industriales. A lo largo del estudio, se comparará este enfoque con otras metodologías existentes en la industria, evaluando su desempeño en términos de tiempos de procesamiento, flexibilidad y adaptabilidad a escenarios reales.

Los sistemas ERP (Enterprise Resource Planning) son herramientas integradas que permiten la automatización y gestión de los procesos empresariales esenciales, tales como finanzas, producción, logística y ventas. Un ERP centraliza la información empresarial, optimizando la toma de decisiones y reduciendo la fragmentación de datos, lo que mejora la eficiencia operativa. Estos sistemas han evolucionado desde los sistemas MRP (Material Requirements Planning) y actualmente constituyen una pieza clave en la digitalización de la industria \cite{ref1}.

\subsection{Automatización con ERP en talleres industriales}

La implementación de ERP en talleres industriales responde a la necesidad de automatizar y estructurar procesos que, de otro modo, se gestionarían de forma manual, lo que puede generar ineficiencias. En los talleres de tecnomecánica, donde se requiere un control riguroso de los repuestos y tiempos de trabajo, un ERP automatizado permite la asignación eficiente de tareas según la disponibilidad de operarios y recursos, reduciendo tiempos de espera y mejorando la productividad. 

En un estudio realizado en 2019 \cite{ref2}, se identificó que el 70.4\% de las quejas de los clientes en una compañía analizada estaban relacionadas con retrasos en la entrega de órdenes de trabajo, mientras que el 29.6\% restante correspondía a problemas de órdenes no confirmadas y reprocesos. Estos problemas, comunes en talleres industriales, pueden ser mitigados mediante un ERP que automatice la distribución de tareas y optimice la administración de recursos, minimizando errores humanos y mejorando la trazabilidad de las órdenes de trabajo.

\subsection{Algoritmos Genéticos (GA) para la Automatización}

La automatización de procesos en entornos industriales ha avanzado con el uso de algoritmos evolutivos, entre los cuales los algoritmos genéticos (GA) han demostrado ser eficaces para resolver problemas de asignación de tareas y planificación de recursos. Estos algoritmos, inspirados en la selección natural y la evolución biológica, trabajan con poblaciones de soluciones potenciales, aplicando operadores como la selección, el cruce (crossover) y la mutación para encontrar soluciones óptimas de manera iterativa \cite{ref3}.

A diferencia de los métodos tradicionales de asignación de tareas, que pueden depender de reglas predefinidas y cálculos deterministas, los algoritmos genéticos ofrecen una solución flexible que se adapta dinámicamente a cambios en el entorno de producción. Gracias a su capacidad de explorar múltiples combinaciones de asignación de tareas y su enfoque probabilístico, los GA pueden reducir tiempos improductivos y mejorar la distribución del trabajo en talleres industriales.

\section{Técnicas de Automatización en Ingeniería Industrial}

La automatización en la ingeniería industrial ha evolucionado con el desarrollo de metodologías como Lean Manufacturing, Six Sigma y la investigación de operaciones, las cuales buscan mejorar la eficiencia mediante la eliminación de desperdicios, la reducción de variabilidad en los procesos y la optimización de los flujos de trabajo \cite{ref4}. Estas estrategias han sido aplicadas con éxito en la manufactura y la gestión industrial, mejorando la toma de decisiones y la productividad.

En este proyecto, el enfoque se basa en la automatización de la asignación de tareas dentro de un ERP mediante inteligencia artificial, permitiendo gestionar dinámicamente la carga de trabajo de los operarios según la demanda y la disponibilidad de recursos. A diferencia de las metodologías tradicionales, que dependen de reglas fijas y modelos matemáticos estrictos, este enfoque permite una mayor flexibilidad y adaptación a entornos productivos cambiantes, mejorando la eficiencia operativa en los talleres industriales.

\section{Antecedentes}

La programación de tareas en entornos industriales ha sido un campo de estudio ampliamente abordado en la literatura debido a la complejidad inherente de problemas como Flow Shop, Job Shop y Flexible Job Shop, todos ellos catalogados como NP-hard \cite{flowshop,flexjobshop}. Esta complejidad ha motivado el desarrollo de metodologías basadas en heurísticas, reglas de despacho y metaheurísticas avanzadas para obtener soluciones eficientes dentro de límites computacionales razonables. En este contexto, los algoritmos genéticos (AG) han demostrado un rendimiento destacado en la optimización de tareas, especialmente en escenarios donde múltiples restricciones interactúan entre sí.

Uno de los trabajos más relevantes es el de Delgado et al. \cite{Delgado2005}, quienes propusieron un algoritmo genético para la programación de tareas en una celda de manufactura y compararon su desempeño con métodos heurísticos tradicionales. Sus resultados mostraron que los AG pueden reducir tiempos improductivos y mejorar la utilización de recursos, destacando su capacidad para explorar soluciones que las reglas de prioridad no alcanzan debido a su naturaleza determinista.

De manera complementaria, Meisel y Prado \cite{meisel2010} desarrollaron un algoritmo genético híbrido que incorpora enfriamiento simulado, demostrando que las metaheurísticas híbridas pueden ofrecer soluciones más estables y de mayor calidad en problemas de tipo Job Shop. Este estudio respalda la idea de que los AG son especialmente adecuados para problemas con múltiples restricciones, como precedencias, máquinas en paralelo y tiempos variables de procesamiento.

En el ámbito de la automatización industrial desde la perspectiva de Industria 4.0, Zhang et al. \cite{industry4} presentan un análisis del estado del arte de algoritmos evolutivos aplicados a la programación en sistemas inteligentes. Su investigación concluye que los AG, junto con técnicas híbridas de optimización, son herramientas clave para la toma de decisiones en entornos altamente dinámicos y automatizados. Esta tendencia refuerza la pertinencia del uso de AG como enfoque moderno para la gestión de recursos en talleres industriales.

En cuanto a la asignación de tareas en problemas de tipo Task Assignment Problem (TAP), Savić et al. \cite{ref10} demostraron que los AG pueden encontrar soluciones óptimas o cercanas al óptimo en espacios de búsqueda complejos, superando métodos deterministas y heurísticos en términos de tiempo de ejecución y calidad de resultados. De forma similar, Wang \cite{ref9} aplicó algoritmos genéticos al diseño automático de maquinaria, mostrando mejoras significativas frente a métodos tradicionales en precisión y tiempos de procesamiento.

Por otra parte, investigaciones sobre heurísticas industriales, como las reglas de despacho SPT, EDD y FCFS, han demostrado ser métodos simples y eficientes para problemas básicos de secuenciación, pero con limitaciones marcadas en problemas con múltiples restricciones o estructuras complejas de precedencia \cite{valletesis,flowshop}. Esta brecha entre simplicidad y capacidad de adaptación es precisamente lo que motiva la comparación entre SPT y algoritmos evolutivos en el presente estudio.

En conjunto, estos antecedentes evidencian que los algoritmos genéticos han sido utilizados con éxito en diversos dominios de la optimización industrial, mientras que las heurísticas tradicionales siguen siendo herramientas fundamentales como métodos comparativos o de referencia. En el caso particular de un taller industrial basado en un sistema ERP, la literatura existente es limitada. Por ello, el presente proyecto aporta un enfoque novedoso al integrar un algoritmo genético dentro de un entorno simulado que replica el flujo real de un taller automotriz, comparándolo directamente con una heurística industrialmente aceptada como SPT. Esta comparación permite evaluar el potencial de los AG para mejorar la asignación de tareas y la eficiencia operativa en un contexto aplicado.

